\scanpage{70}
\mspage{63}
% archive sheet number: 69
\uline{Lemme} La catégorie
\diagram{
$\unsure{\Psi|\Upsilon}={}$
\begin{tikzcd}[ampersand replacement=\&,row sep=small]
\&1\\
0\arrow[ur]\arrow[dr]\&\\
\&2
\end{tikzcd}
}
est admissible.

Ceci va résulter du

\uline{Cor.} Soit
\diagram{
$X_\bullet={}$
\begin{tikzcd}[ampersand replacement=\&,row sep=small]
\&X_1\arrow[dl,"\alpha_1"']\\
X_0\&\\
\&X_2\arrow[ul,"\alpha_2"]
\end{tikzcd}
}
dans $M(\Psi)$. Pour $X_\bullet$ soit fibrant, il f. et s. que

a) $X_0,X_1,X_2$ soient fibrants et

b) $\alpha_1,\alpha_2\in F$.
% margin note written slanted
\margin{\uline{NB} Il y a une variante évidente \unsure{caractérisant} les $X$ tels que $(X\unsure{-|\to}e)$ soit dans $\mathrm{TF}(\Psi)$, \unsure{comme ceux où|comme pour} $X_0,X_1,X_2$ sont \deleted{\unsure{f}}\unsure{fibrants} \uline{triviaux}, et $\alpha_1,\alpha_2\in\mathrm{TF}$.}

La nécessité de a) est \inserted{cas particulier} du cor. 2 page 48, celle de b) est cas part. du cor. 1 p 50.

Pour la suffisance, \inserted{\unsure{(dernier exercice préliminaire pour le th.)|(exercice préliminaire pour le th.)}} on \deleted{\unsure{doit|soit}} montre que dans le diagramme comm. à traits pleins
\diagram{
\begin{tikzcd}[ampersand replacement=\&,column sep=small,row sep=large]
A_{\unsure{1|0|\rho}}\arrow[rr]\arrow[dr,"\mathrm{TC}"']\arrow[dd] \& \& A_0\arrow[dr,"\mathrm{TC}"]\arrow[dd] \& \& A_2\arrow[ll]\arrow[dr,"\mathrm{TC}"]\arrow[dd] \& \\
\& B_1\arrow[rr]\arrow[dl,dashed,"h_1"'] \& \& B_0\arrow[dl,dashed,"h_0"'] \& \& B_2\arrow[ll]\arrow[dl,dashed,"h_2"] \\
X_1\arrow[rr,"\alpha_1"'] \& \& X_0 \& \& X_2\arrow[ll,"\alpha_2"] \&
\end{tikzcd},
}
avec trois flèches marquées $\mathrm{TC}$, on peut trouver $h_1,h_0,h_2$ laissant le diagr. comm. (trois triangles et deux carrés comm. rajoutés). Pour ceci, on commence par choisir $h_0$ de façon à rendre le triangle vertical médian comm. --- c'est possible car $X_0$ est fibrant, $A_0\unsure{-|\to}B_0$ dans $\mathrm{TC}$.
